% Corrected French for printed p.131.
% Corrections: schémas -> diagrammes; extra opening parenthesis removed;
% momorphisme -> morphisme; formé -> formée.
\SourcePage{131}

\LayerTarget{tohoku:def:additive-diagram-relations:p2}
\noindent $R_{ij}$, on dira que $D$ est $R$-commutatif si tous les $D(L)$
$(L\in R)$ sont nuls. Appelons \textit{schéma de diagramme avec relations de
commutation} un couple $(\mathrm{S},R)=\Sigma$ formé d'un schéma de diagramme et d'un
ensemble $R$ comme ci-dessus. Pour toute catégorie additive $\mathbf{C}$, on peut alors
considérer la sous-catégorie $\mathbf{C}^\Sigma$ de $\mathbf{C}^S$ formée des diagrammes
$R$-commutatifs.

\LayerTarget{tohoku:prop:1.6.1}
\textsc{Proposition 1.6.1.} \textit{Soient $\Sigma$ un schéma de
diagramme avec relations de commutation, $\mathbf{C}$ une catégorie additive. Alors la
catégorie $\mathbf{C}^\Sigma$ est une catégorie additive, et si $\mathbf{C}$ est une catégorie
avec produits directs infinis (resp. sommes directes infinies) il en est de
même de $\mathbf{C}^\Sigma$. De plus, si $\mathbf{C}$ satisfait à un des axiomes \textup{AB 1)} à \textup{AB 6)}
ou des axiomes duals \textup{AB 3*)} à \textup{AB 6*)}, il en est de même de $\mathbf{C}^\Sigma$.}

\LayerTarget{tohoku:prop:1.6.1:proof}
D'ailleurs, si $D,D'\in \mathbf{C}^\Sigma$, et si $u$ est un morphisme de $D$ dans
$D'$, alors son noyau (resp. conoyau, resp. image, resp. coimage) est le
diagramme formé par les noyaux (resp. \ldots) des composantes $u_i$, les
morphismes dans ce diagramme (correspondants aux flèches du schéma) se
déduisant de ceux de $D$ (resp. $D'\ldots$) de la façon évidente par
restriction (resp. passage au quotient). On interprète de façon analogue la
somme directe ou le produit direct d'une famille de diagrammes. Les sous-trucs
$D'$ d'un diagramme $D$ s'identifient aux familles $(D'(i))$ de sous-trucs des
$D(i)$ telles que pour toute flèche $\varphi$ d'origine $i$ et extrémité $j$,
on ait $D(\varphi).\,D'(i)\subset D'(j)$ ; alors $D'(\varphi)$ se définit comme
le morphisme $D'(i)\to D'(j)$ défini par $D(\varphi)$. Les trucs
quotients de $D$ se déterminent de façon duale.

\LayerTarget{tohoku:def:dual-diagram-scheme}
Si $\mathrm{S}$ est un schéma de diagrammes, on appelle \textit{schéma dual} et on
désigne par $\mathrm{S}^0$ le schéma ayant mêmes sommets et même ensemble de flèches
que $\mathrm{S}$, mais l'origine et l'extrémité des flèches de $\mathrm{S}$ étant intervertis. Si
on donne de plus un ensemble $R$ de relations de commutation pour $\mathrm{S}$, on
gardera le même ensemble pour $\mathrm{S}^0$. Avec cette convention, pour une
catégorie additive $\mathbf{C}$, la catégorie duale de $\mathbf{C}^\Sigma$ s'identifie à
$(\mathbf{C}^\circ)^{\Sigma^\circ}$.

\LayerTarget{tohoku:def:canonical-diagram-extension}
Soient $\mathbf{C},\mathbf{C}'$ deux catégories additives, $\Sigma$ un schéma de diagramme avec
relations de commutation. Pour tout foncteur $F$ de $\mathbf{C}$ dans $\mathbf{C}'$, on définit
de façon évidente le foncteur $F^\Sigma$ de $\mathbf{C}^\Sigma$ dans $\mathbf{C}'^\Sigma$, dit
\textit{prolongement canonique de $F$ aux diagrammes}. $F^\Sigma$ se comporte
formellement comme un foncteur par rapport à l'argument $F$, en particulier un
morphisme fonctoriel $F\to F'$ définit un homomorphisme fonctoriel
$F^\Sigma\to F'^\Sigma$. Notons enfin qu'on a pour un foncteur
composé : $(GF)^\Sigma=G^\Sigma F^\Sigma$, et que les propriétés d'exactitude d'un foncteur se conservent par
prolongement à une classe de diagrammes.

\LayerTarget{tohoku:section:1.7}
\textbf{1.7. Exemples de catégories définies par des schémas de
diagrammes.}
\LayerTarget{tohoku:example:1.7:a}
a) Prenons $I$ réduit à un élément, et l'ensemble des flèches vide.
Alors les relations de commutation sont de la forme $n_i e=0$, donc peuvent se
réduire à une unique relation $ne=0$. Alors $\mathbf{C}^\Sigma$ est la sous-catégorie de
$\mathbf{C}$ formée des objets annulés par l'entier $n$. Si $n=0$, c'est $\mathbf{C}$ lui-même.

\LayerTarget{tohoku:example:1.7:b:p1}
b) Prenons $I$ quelconque, pas de flèches, pas de relations de
commutation. Alors $\mathbf{C}^\Sigma$ s'identifie à la catégorie produit $\mathbf{C}^I$. Si on
suppose données des

\TohokuContinuation{132}
